\documentclass[a4paper,11pt]{article}
\usepackage[T2A]{fontenc}
\usepackage[utf8]{inputenc}
\usepackage[russian]{babel}
\usepackage{amsmath,amssymb,mathtools}
\usepackage{icomma}
\usepackage[a4paper,left=23mm,right=23mm,top=22mm,bottom=24mm]{geometry}
\usepackage{microtype}
\usepackage{tikz}
\usetikzlibrary{arrows.meta,positioning,shapes.geometric,calc}
\usepackage{pgfplots}
\pgfplotsset{compat=1.18}
\usepackage{booktabs,tabularx,array}
\usepackage{enumitem}
\usepackage{fancyhdr}
\usepackage{setspace}
\usepackage{xcolor}


\definecolor{accent}{HTML}{008080}
\definecolor{darkaccent}{HTML}{004D4D}
\definecolor{textgray}{HTML}{333333}
\definecolor{pale}{HTML}{788B8B}
\color{textgray}
\setstretch{1.27}
\setlength{\parindent}{0pt}
\setlength{\parskip}{0.24em}
\setlist[itemize]{leftmargin=6mm,itemsep=0.08em,topsep=0.10em}
\setlist[enumerate]{leftmargin=7mm,itemsep=0.18em,topsep=0.12em}
\setlength{\headheight}{15pt}
\pagestyle{fancy}
\fancyhf{}
\fancyhead[L]{\footnotesize\textcolor{pale}{Задачи прикладного характера}}
\fancyhead[R]{\footnotesize\textcolor{pale}{Анна Трапезникова}}
\fancyfoot[C]{\small\textcolor{pale}{\thepage}}
\renewcommand{\headrulewidth}{0.25pt}
\renewcommand{\footrulewidth}{0pt}
\renewcommand{\arraystretch}{1.24}
\newcommand{\key}[1]{\textcolor{darkaccent}{\textbf{#1}}}
\newcommand{\ex}[1]{\par\smallskip\textbf{\textcolor{accent}{#1}}\par}
\newcommand{\CheckItem}{\(\square\)}

\begin{document}
\begin{titlepage}
\thispagestyle{empty}
\centering
\vspace*{18mm}
{\large\color{darkaccent} МАТЕМАТИКА\par}
\vspace{19mm}
{\LARGE Чек-лист\par}
\vspace{11mm}
{\Large\color{accent} Задачи прикладного характера\par}
{\Large\color{accent} в профильном ЕГЭ\par}
\vspace{16mm}
{\normalsize Формулы, единицы измерения, линейные\\
и квадратные уравнения\par}
\vfill
{\large Анна Трапезникова\par}
\vspace{15mm}
{26 октября 2026 г.\par}
\vspace{16mm}
{\small Лёвин Артём Александрович\par}
{\small эксклюзивно для Анны Трапезниковой\par}
\vspace{18mm}
\begin{tikzpicture}
\draw[accent,line width=0.75pt]
(0,0) .. controls (1.0,0.44) and (2.1,-0.42) .. (3.2,0)
      .. controls (4.3,0.47) and (5.45,-0.43) .. (6.6,0)
      .. controls (7.7,0.40) and (8.7,-0.31) .. (9.4,0);
\end{tikzpicture}
\vspace{8mm}
\end{titlepage}

\section*{1. Что требуется уметь}
В задачах прикладного характера готовая формула связывает несколько величин.
Главная работа --- понять условие, выбрать нужную величину и грамотно выполнить подстановку.

\begin{itemize}
\item[\CheckItem] Выделять \key{известные величины}, неизвестную и их единицы измерения.
\item[\CheckItem] Переводить словесные условия в равенства и неравенства.
\item[\CheckItem] Понимать разницу между исходным значением, новым значением и изменением.
\item[\CheckItem] Приводить единицы к одной системе и вычислять степени десяти.
\item[\CheckItem] Решать линейные и квадратные уравнения, проверяя допустимость ответа.
\item[\CheckItem] Выбирать наибольшее, наименьшее или физически осмысленное решение.
\end{itemize}

\subsection*{Универсальный маршрут}
\key{Прочитайте условие} $\longrightarrow$ \key{выпишите данные} $\longrightarrow$
\key{найдите связь} $\longrightarrow$ \key{подставьте числа} $\longrightarrow$
\key{решите} $\longrightarrow$ \key{проверьте смысл ответа}.

\section*{2. Как переводить текст на язык формул}

\begin{tabularx}{\linewidth}{@{}p{.35\linewidth}X@{}}
\toprule
Формулировка & Математическая запись\\
\midrule
Увеличилась на \(a\) & \(x_{\text{нов}}-x_{\text{стар}}=a\)\\
Уменьшилась на \(a\) & \(x_{\text{стар}}-x_{\text{нов}}=a\)\\
Не меньше \(a\) & \(x\ge a\)\\
Не больше \(a\) & \(x\le a\)\\
Наибольшее допустимое значение & Найти все допустимые значения и выбрать максимальное\\
Минимальная скорость & Найти порог, при котором условие выполняется\\
\bottomrule
\end{tabularx}

\textcolor{accent}{\textbf{Обозначение функции.}}
Запись \(L(t)\) означает, что длина зависит от температуры \(t\).
При подстановке конкретного \(t\) можно работать с буквой \(L\), сохраняя смысл формулы.

\textcolor{accent}{\textbf{Важное уточнение.}}
Слова «не меньше» и «не больше» определяют знак неравенства.
Для монотонной линейной зависимости граничное значение удобно находить из
равенства, но затем необходимо проверить направление неравенства.

\section*{3. Единицы и степени десяти}
\[
1\text{ мм}=10^{-3}\text{ м},\qquad
4{,}5\text{ мм}=4{,}5\cdot10^{-3}\text{ м}
=45\cdot10^{-4}\text{ м}.
\]
\[
1{,}2\cdot10^{-5}=12\cdot10^{-6},\qquad
\frac{10^a}{10^b}=10^{a-b}.
\]
\clearpage
\section*{4. Тепловое расширение: изменение длины}

При температуре \(0^\circ\mathrm{C}\) длина рельса составляет \(L_0=10\) м.
При нагревании она меняется по формуле
\[
L(t)=L_0(1+\alpha t), \qquad
\alpha=1{,}2\cdot10^{-5}\;{}^\circ\mathrm{C}^{-1}.
\]
Найдём температуру, при которой рельс удлинится на \(4{,}5\) мм.

\ex{Шаг 1. Выразите изменение длины.}
«Удлинился на \(4{,}5\) мм» означает
\[
L-L_0=\Delta L=4{,}5\cdot10^{-3}\text{ м}.
\]

\ex{Шаг 2. Упростите исходную формулу.}
\[
L=L_0+L_0\alpha t
\quad\Longrightarrow\quad
\Delta L=L_0\alpha t.
\]

\ex{Шаг 3. Выразите температуру и подставьте числа.}
\[
t=\frac{\Delta L}{L_0\alpha}
 =\frac{4{,}5\cdot10^{-3}}
 {10\cdot1{,}2\cdot10^{-5}}
 =\frac{4{,}5}{1{,}2}\cdot10^{1}
 =37{,}5.
\]
\key{Ответ:} \(37{,}5^\circ\mathrm C\).

\textbf{Проверка:}
\(10\cdot1{,}2\cdot10^{-5}\cdot37{,}5=0{,}0045\) м,
то есть ровно \(4{,}5\) мм.

\section*{5. Прибыль предприятия: линейная зависимость}
Месячная прибыль описана формулой
\[
\Pi(q)=q(p-v)-F,
\]
где \(q\) --- количество изделий,
\(p\) --- цена одного изделия,
\(v\) --- затраты на изготовление одного изделия,
\(F\) --- постоянные расходы.

Пусть \(p=600\) руб., \(v=300\) руб., \(F=700\,000\) руб.
Найдём минимальный объём производства, при котором прибыль
\textbf{не меньше \(500\,000\) руб.}
\[
q(600-300)-700\,000\ge500\,000,
\]
\[
300q\ge1\,200\,000
\quad\Longrightarrow\quad q\ge4000.
\]
\key{Ответ:} \(4000\) изделий.

\textbf{Почему границу можно найти равенством?}
Здесь каждое дополнительное изделие увеличивает прибыль на \(300\) руб.,
поэтому при \(q=4000\) достигается нужный порог.
Если расчёт даёт дробное число изделий, результат для минимального
целого объёма следует \textbf{округлить вверх}.

\clearpage
\section*{6. Колодец: косвенно заданное время}

Пока камень падает, пройденное расстояние определяется формулой
\[
h=5t^2,
\]
где \(h\) измеряется в метрах, \(t\) --- в секундах.
До дождя время падения было \(t_1=0{,}6\) с,
а после дождя уменьшилось на \(0{,}2\) с.

\begin{center}
\begin{tikzpicture}[scale=.85,>=Latex]
\foreach \x/\name/\yy in {0/До дождя/0.35,6/После дождя/1.05}{
  \draw[darkaccent,thick] (\x,2.8)--(\x,.1) (\x+2.0,2.8)--(\x+2,.1);
  \draw[accent,thick] (\x,\yy)--(\x+2,\yy);
  \fill[accent!13] (\x+.04,.12) rectangle (\x+1.96,\yy-.02);
  \node[above] at (\x+1,3.0) {\small\name};
  \draw[<->,pale] (\x+1.45,2.7)--(\x+1.45,\yy+.05);
}
\node[below] at (1,.02) {\small \(t_1=0{,}6\text{ с}\)};
\node[below] at (7,.02) {\small \(t_2=0{,}4\text{ с}\)};
\end{tikzpicture}
\end{center}

\ex{Шаг 1. Найдите время после дождя.}
\[
t_2=0{,}6-0{,}2=0{,}4\text{ с}.
\]

\ex{Шаг 2. Вычислите сокращение глубины до поверхности воды.}
\[
\Delta h=h_1-h_2=5t_1^2-5t_2^2
=5(0{,}6^2-0{,}4^2).
\]
Используем формулу разности квадратов:
\[
\Delta h=5(0{,}6-0{,}4)(0{,}6+0{,}4)
=5\cdot0{,}2\cdot1=1\text{ м}.
\]
\key{Ответ:} уровень воды поднялся на \(1\) м.

\textbf{Проверьте смысл:} меньшая продолжительность падения означает,
что поверхность воды стала ближе к верхнему краю колодца.

\section*{7. Две формулы: спрос и выручка}

Пусть объём спроса \(q=60-5p\), а выручка \(R=qp\).
Здесь цена \(p\) измеряется в тыс. руб., выручка \(R\) --- в тыс. руб.
Найдём \textbf{наибольшую цену}, при которой \(R\ge160\).
\[
R=p(60-5p)\ge160,
\]
\[
-5p^2+60p-160\ge0
\quad\Longleftrightarrow\quad
p^2-12p+32\le0.
\]
Корни граничного уравнения:
\[
(p-4)(p-8)=0
\quad\Longrightarrow\quad p_1=4,\quad p_2=8.
\]
Парабола \(p^2-12p+32\) направлена ветвями вверх;
неравенство выполняется между корнями:
\[
4\le p\le8.
\]
\key{Ответ:} максимальная цена \(8\) тыс. руб.

\textbf{Проверка:} \(q=60-5\cdot8=20\ge0\) и
\(R=8\cdot20=160\) тыс. руб.

\section*{8. Ведро на верёвке: физический смысл корней}

Сила давления воды на дно в верхней точке движения по окружности
выражается формулой
\[
P=m\left(\frac{v^2}{L}-g\right),
\]
где \(m>0\) --- масса воды, \(L>0\) --- длина верёвки,
\(v\) --- скорость, \(g=10\text{ м/с}^2\).
Условие удержания воды в верхней точке: \(P\ge0\).

Поскольку \(m>0\), получаем
\[
\frac{v^2}{L}-g\ge0
\quad\Longrightarrow\quad v^2\ge gL.
\]
При \(L=0{,}625\) м критическое значение скорости:
\[
v_{\min}=\sqrt{gL}=\sqrt{10\cdot0{,}625}
=\sqrt{6{,}25}=2{,}5\text{ м/с}.
\]
\key{Ответ:} \(2{,}5\) м/с.

Значение \(v=-2{,}5\) алгебраически соответствует уравнению
для квадрата скорости, однако в условии требуется \textbf{модуль скорости},
поэтому выбираем неотрицательное значение. Масса \(m=0\) исключается
по смыслу задачи: вода в ведре есть.

\section*{9. Параболическая траектория: выбираем крайний корень}

Высота летящего камня \(y\) связана с горизонтальным расстоянием \(x\)
формулой
\[
y=-\frac{x^2}{60}+\frac76x.
\]
Для стены высотой \(9\) м необходим запас не менее \(1\) м,
поэтому в точке стены \(y\ge9+1=10\text{ м}\).
Подставим высоту в условие:
\[
-\frac{x^2}{60}+\frac76x\ge10
\quad\Longleftrightarrow\quad
x^2-70x+600\le0.
\]
\[
(x-10)(x-60)\le0
\quad\Longrightarrow\quad 10\le x\le60.
\]
\key{Ответ:} максимальное расстояние \(60\) м.

\clearpage
\section*{Графическая проверка траектории}
На графике показаны точки, в которых высота камня равна требуемым
\(10\) м. Между ними траектория проходит выше этой отметки.

\begin{center}
\begin{tikzpicture}
\begin{axis}[
 width=.85\linewidth,height=4.7cm,
 xmin=0,xmax=70,ymin=0,ymax=23,
 axis lines=left,
 xlabel={\(x\), м},ylabel={\(y\), м},
 xtick={0,10,35,60,70},ytick={0,9,10,20},
 tick label style={font=\footnotesize},
 label style={font=\small},
 grid=major,grid style={black!8},
 every axis plot/.append style={line width=1.15pt},
 clip=false
]
\addplot[accent,domain=0:70,samples=220]
 {-x^2/60+7*x/6};
\addplot[dashed,black!45,domain=0:70] {10};
\addplot[only marks,mark=*,mark size=1.9pt,darkaccent]
 coordinates {(10,10) (60,10)};
\node[anchor=south west,font=\scriptsize] at (axis cs:13,10) {запас \(1\) м};
\node[anchor=north east,font=\scriptsize] at (axis cs:60,10) {крайнее положение};
\end{axis}
\end{tikzpicture}
\end{center}

\smallskip
\begin{center}
\begin{tabularx}{.85\linewidth}{@{}p{.19\linewidth}p{.34\linewidth}X@{}}
\toprule
Расстояние \(x\) & Высота \(y(x)\) & Условие \(y\ge10\)\\
\midrule
\(10\) м & \(10\) м & Выполнено\\
\(35\) м & \(\dfrac{245}{12}\) м & Выполнено\\
\(60\) м & \(10\) м & Выполнено\\
\(65\) м & \(\dfrac{65}{12}\) м & Не выполнено\\
\bottomrule
\end{tabularx}
\end{center}
\textbf{Проверка:} при \(x=60\) высота камня равна \(10\) м;
при большем расстоянии камень уже проходит ниже требуемой отметки.

\clearpage
\section*{10. Алгоритм решения прикладной задачи}
\begin{center}
\begin{tikzpicture}[
>=Latex,
node distance=11mm,
start/.style={ellipse,draw=darkaccent,very thick,minimum width=35mm,minimum height=10mm,align=center,font=\small},
action/.style={rounded corners=2pt,draw=accent,minimum width=80mm,minimum height=13mm,align=center,text width=74mm,font=\small},
dec/.style={diamond,draw=darkaccent,aspect=2.25,inner sep=2.5pt,align=center,text width=35mm,font=\small},
arr/.style={->,very thick,accent},
corr/.style={rounded corners=2pt,draw=pale,minimum width=34mm,minimum height=12mm,align=center,text width=32mm,font=\small}
]
\node[start] (a) {Условие};
\node[action,below=of a] (b) {Выпишите данные, искомую величину, единицы и ограничения};
\node[action,below=of b] (c) {Выберите формулу, выразите неизвестную и подставьте числа};
\node[action,below=of c] (d) {Решите уравнение или неравенство};
\node[dec,below=14mm of d] (e) {Все значения подходят по смыслу?};
\node[start,below=15mm of e] (f) {Запишите ответ с единицами};
\node[corr,right=26mm of e] (g) {Отберите допустимые значения};
\draw[arr] (a)--(b);
\draw[arr] (b)--(c);
\draw[arr] (c)--(d);
\draw[arr] (d)--(e);
\draw[arr] (e)--node[left,font=\footnotesize]{да}(f);
\draw[arr] (e)--node[above,font=\footnotesize]{нет}(g);
\draw[arr] (g.south) |- ($(f.north)+(0,0.0)$);
\end{tikzpicture}
\end{center}
\vfill
\begin{center}
{\small\textcolor{pale}{Один неверно выбранный корень может испортить весь верный расчёт.}}
\end{center}

\clearpage
\section*{11. Частые ошибки}
\begin{enumerate}
\item \textbf{Подмена изменения конечным значением.}
   Если длина выросла на \(4\) мм, то \(\Delta L=4\) мм,
   а полная длина равна \(L_0+\Delta L\).
\item \textbf{Смешивание метров и миллиметров.}
   Перед подстановкой переводите все длины в согласованные единицы.
\item \textbf{Игнорирование слов «не менее» и «не более».}
   Записывайте соответствующий знак; равенство используется для поиска границы.
\item \textbf{Неверное время после изменения.}
   Если время сократилось на \(0{,}2\) с, новое время равно \(t_1-0{,}2\).
\item \textbf{Потеря второго корня.}
   Квадратное уравнение может иметь два допустимых ответа;
   условие «наибольший» указывает, какой выбрать.
\item \textbf{Деление на нуль или недопустимый физический случай.}
   Учитывайте положительность массы, длины, объёма и другие ограничения.
\end{enumerate}

\section*{Итоговая памятка}
\begin{tabularx}{\linewidth}{@{}p{.39\linewidth}X@{}}
\toprule
Сигнал в условии & Следующее действие\\
\midrule
«Увеличится на \(\ldots\)» & Записать разность нового и старого значения.\\
Два разных набора единиц & Перевести к одной системе.\\
Формула с \(L(t)\), \(q(p)\) & Прочитать зависимость, найти нужную переменную.\\
«Не меньше» / «не больше» & Составить неравенство и проверить границу.\\
Две формулы с общей величиной & Использовать подстановку.\\
Квадратное уравнение & Найти корни и отобрать по условию.\\
Сюжет о движении & Сделать схему, проверить физический смысл.\\
\bottomrule
\end{tabularx}

\clearpage
\section*{Тренировочный блок}
\textcolor{accent}{\textbf{Путь А · 10 задач}}
\par
Решайте последовательно. Выписывайте исходные данные и ограничения,
где они необходимы. В ответ записывайте требуемые единицы.

\subsection*{Базовый уровень}
\begin{enumerate}[label=\textbf{\arabic*.}]
\item Брус длиной \(8\) м удлинился на \(3{,}6\) мм.
      Найдите новую длину бруса в метрах.
\item Представьте число \(0{,}000072\) в стандартном виде
      \(a\cdot10^n\), где \(1\le a<10\) и \(n\) --- целое число.
\item Длина рельса при \(0^\circ\mathrm C\) равна \(12\) м.
      Зависимость длины от температуры:
      \(L(t)=L_0(1+\alpha t)\), где
      \(\alpha=1{,}5\cdot10^{-5}\;{}^\circ\mathrm C^{-1}\).
      При какой температуре рельс удлинится на \(9\) мм?
\item Длина металлического стержня при \(0^\circ\mathrm C\) равна \(8\) м,
      \(\alpha=1{,}25\cdot10^{-5}\;{}^\circ\mathrm C^{-1}\).
      Найдите удлинение стержня в миллиметрах при \(t=40^\circ\mathrm C\).
      Используйте формулу \(L(t)=L_0(1+\alpha t)\).
\end{enumerate}

\subsection*{Средний уровень}
\begin{enumerate}[label=\textbf{\arabic*.},start=5]
\item Прибыль предприятия задаётся формулой
      \(\Pi(q)=q(p-v)-F\).
      Цена изделия \(p=800\) руб.,
      переменные затраты \(v=500\) руб.,
      постоянные расходы \(F=420\,000\) руб.
      Найдите минимальное целое \(q\), при котором
      прибыль не меньше \(180\,000\) руб.
\item Цена изделия \(700\) руб., переменные затраты \(460\) руб.,
      постоянные расходы \(36\,000\) руб.
      По формуле \(\Pi(q)=q(p-v)-F\) найдите минимальное
      \textbf{целое} число изделий, необходимое для прибыли
      не менее \(120\,100\) руб.
\item Глубина, пройденная камнем при падении, задаётся формулой
      \(h=5t^2\). До дождя камень падал до воды \(0{,}9\) с.
      После дождя время падения сократилось на \(0{,}3\) с.
      На сколько метров поднялся уровень воды?
\end{enumerate}

\subsection*{Уровень выше среднего}
\begin{enumerate}[label=\textbf{\arabic*.},start=8]
\item Объём спроса описан формулой \(q=80-4p\),
      выручка --- формулой \(R=qp\), где \(p\) в тыс. руб.,
      \(R\) в тыс. руб.
      Найдите наибольшую цену \(p\), при которой \(R\ge300\).
\item В верхней точке окружности сила давления воды в ведре
      равна \(P=m\bigl(v^2/L-g\bigr)\).
      При \(L=0{,}9\) м и \(g=10\text{ м/с}^2\)
      найдите минимальную скорость \(v\) (в м/с),
      при которой \(P\ge0\). Считайте \(m>0\).
\item Камень летит по траектории
      \(y=-\dfrac{x^2}{20}+2x\), где \(x\) и \(y\) измеряются
      в метрах. На каком \textbf{наибольшем} расстоянии
      от стены высотой \(14\) м можно расположить метательную
      машину, чтобы камень проходил над стеной
      с запасом не менее \(1\) м?
\end{enumerate}

\clearpage
\section*{Ответы}
\begin{center}
\begin{tabularx}{.91\linewidth}{@{}p{.15\linewidth}X@{}}
\toprule
№ & Ответ\\
\midrule
1 & \(8{,}0036\) м\\
2 & \(7{,}2\cdot10^{-5}\)\\
3 & \(50^\circ\mathrm C\)\\
4 & \(4\) мм\\
5 & \(2000\) изделий\\
6 & \(651\) изделие\\
7 & \(2{,}25\) м\\
8 & \(15\) тыс. руб.\\
9 & \(3\) м/с\\
10 & \(30\) м\\
\bottomrule
\end{tabularx}
\end{center}

\vspace{1cm}
\section*{Самопроверка результата}
После решения каждой задачи убедитесь, что:
\begin{itemize}
\item ответ имеет нужные единицы;
\item все использованные величины относятся к одной системе единиц;
\item выбранный корень соответствует словам «наибольший» или «наименьший»;
\item числовой результат согласуется с реальным смыслом задачи.
\end{itemize}
\vfill
\begin{center}
\small\textcolor{pale}{Лёвин Артём Александрович эксклюзивно для Анны Трапезниковой}
\end{center}
\end{document}
